\documentclass{../note}

\begin{document}

\begin{zadanie}{1}
Funkcja $f : \mathbb{R} \to \mathbb{R}$ jest różniczkowalna oraz spełnia $\lim_{x\to 0} \frac{f'(x)}{x^3} = 1$. Oblicz granicę
\[ \lim_{x\to 0} \frac{f(x) - f(\sin x)}{x^3}. \]
\end{zadanie}
\begin{rozwiazanie}
Z założenia wynika, że $f'(x) = x^3 + o(x^3)$. Z definicji pochodnej $f$:
\[\frac{f(x+h) - f(x)}{h} \xrightarrow{h\to0} f'(x)\]
możemy wywnioskować, że $f(x + h) - f(x) = f'(x)h + o(h)$. $\sin x = x - \frac{x^3}{6} + o(x^3)$, więc
\[f(\sin x) - f(x) = f\left(x - \frac{x^3}{6} + o(x^3)\right) - f(x) = f'(x)\left(- \frac{x^3}{6} + o(x^3)\right) + o\left(- \frac{x^3}{6} + o(x^3)\right) =\]\[ \left(x^3 + o(x^3)\right)\left(- \frac{x^3}{6} + o(x^3)\right) + o(x^3) = o(x^3)\]
Zatem $\lim_{x\to 0} \frac{f(x) - f(\sin x)}{x^3} = 0$.
\end{rozwiazanie}

\begin{zadanie}{2}
W wyniku obrotu trójkąta prostokątnego o bokach $a, b, c$, spełniających warunek $a^2 + b^2 = c^2$, wokół jego przeciwprostokątnej powstaje bryła o objętości $V$.
Wyznacz wymiary trójkąta prostokątnego o obwodzie
\[ a + b + c = 1, \]
dla którego objętość $V$ jest maksymalna.
\end{zadanie}
\begin{rozwiazanie}
Niech $h$ to wysokość opuszczona na $c$ oraz niech $p$ to odcinek, który tworzy trójkąt z $h$ i $a$. Skoro $1 - a - c = b = \sqrt{c^2 - a^2}$. Podnosząc tę równość do kwadratu i porządkując wyrazy dostajemy $c = \frac{2a - 2a^2 - 1}{2a - 2} = \frac{1}{2-2a} - a$. Objętość bryły to suma objętości dwóch stożków:
\[\frac{\pi}{3} ph^2 + \frac{\pi}{3} (c-p)h^2 = \frac{\pi}{3} ch^2\]
 Z podobieństwa trójkątów $\frac{h}{a} = \frac{b}{c}$, czyli $h = \frac{ab}{c}$. Do tego $b = 1 - a - c = 1 - \frac{1}{2-2a} = \frac{1-2a}{2-2a}$
\end{rozwiazanie}

\end{document}